Ex\+A\+CT is the framework for data Extraction and Analysis for the Cherenkov Telescope. It creates a library called lib\+Ex\+A\+CT for access to the files created.

\subsection*{Pre-\/\+Requisites}

R\+O\+OT Analysis. The installation has only been tested with versions 6 and above. Event\+Builder (optional). The event builder provided from E\+U\+S\+O-\/\+S\+P\+B2/\+CT repository. This is used to generate the input files for Ex\+A\+CT.

\subsection*{Installation}

The installation of Ex\+A\+CT is done with Make. In the main directory /exact/


\begin{DoxyCode}
make
\end{DoxyCode}
 Will run the {\itshape Makefile} and compile all necessary files and generate the library {\itshape lib\+Ex\+Act.\+so}. This library will be copied to the R\+O\+OT shared libraries for later use.

The executable {\itshape Ex\+A\+CT} will be created and placed in the \+\_\+./exact/\+\_\+ directory

The environment variable {\itshape L\+D\+\_\+\+L\+I\+B\+R\+A\+R\+Y\+\_\+\+P\+A\+TH} must be appended with the path pointing to the dictionaries generated after compilation. Namely, \+\_\+./exact/dict\+\_\+


\begin{DoxyCode}
#.bashrc

LD\_LIBRARY\_PATH=$\{LD\_LIBRARY\_PATH\}:[/path/to/]exact/dict
\end{DoxyCode}


\subsection*{Usage}

To run {\itshape Ex\+A\+CT}\+:


\begin{DoxyCode}
./ExACT -c ./ExACT.cfg -df [path/to/]InputFile -o [path/to/OutputDirectory] -uw [true/false]
\end{DoxyCode}


This provides the compiled code with the needed information for running.

\+\_\+./\+Ex\+A\+CT.cfg\+\_\+ is the configuration file with information about the extraction and analysis to be performed to the data. A copy of the structure is provided in the repository. Active parameters are signaled with an {\bfseries $\ast$} and if a parameter is not to be used it can be \char`\"{}turned off\char`\"{} by adding a \+\_\+\+\_\+\+::\+\_\+\+\_\+ in front. The file also provides a description of each of the parameters.

{\itshape \mbox{[}path/to/\mbox{]}Input\+File} It is the Event\+Builder generated file. The raw data is merged with the trigger board information by the E\+Vent\+Builder and outputs a file where the traces for each pixel are saved. This file is the one used for extraction and analysis.

{\itshape \mbox{[}path/to/\+Output\+Directory\mbox{]}} The directory where the output file should be saved.

{\itshape \mbox{[}true/false\mbox{]}} Whether to use in-\/flight analysis \mbox{[}false\mbox{]} or not \mbox{[}true\mbox{]}

After extraction Ex\+A\+CT generates an output R\+O\+OT file in the same directory as the input file named {\itshape Input\+File\+\_\+\+Extracted.\+root}. This file is populated with a tree called \char`\"{}t1\char`\"{}. Inside this tree each branch corresponds to a different pixel in the camera. In the case of the CT spu it means that the output file contains 512 branches, each with an entry per event. The branches are of data-\/type Extracted\+Data, which is defined in the {\itshape \hyperlink{ExtractedData_8h_source}{Extracted\+Data.\+h}} and {\itshape Extracted\+Data.\+cpp}.

In case it is specified in the configuration file, the naming convention of the extracted file can change if F\+L\+I\+G\+H\+T\+M\+O\+DE is enabled. In that case, the output file follows the naming convention for the download of files through the G\+CC.

\subsubsection*{Reading Output File}

A library called lib\+Ex\+A\+C\+T.\+so contains information on how to access the output R\+O\+OT files generated by Ex\+A\+CT. A sample macro is provided in \+\_\+/exact/macros/\+Read\+Extracted.\+cpp\+\_\+. In order to be able to use the classes defined by Ex\+A\+CT, the following macro should be included in the header of the code to be used\+:


\begin{DoxyCode}
\{c++\}
R\_\_LOAD\_LIBRARY(libExACT.so)
\end{DoxyCode}


Alternatively it is possible to load the library interactively in the R\+O\+OT C\+LI using \+\_\+.\+L\+\_\+


\begin{DoxyCode}
\{c++\}
root [0] .L [path/to/exact/libExACT.so]
root [1] [Code]
\end{DoxyCode}


\subsection*{Contributing}

Please use a branch under dev to make any modifications and tests before pushing anything to the main branch.

Please follow the established conventions below\+:

For C\+L\+A\+S\+S\+ES\+:

Variables\+:
\begin{DoxyItemize}
\item All public variables should start with capital letter I

e.\+g. int Inumber\+Of\+Photons
\item Only camel casing should be used. No underscores for variable names

e.\+g. this\+Is\+Camel\+Casing this\+\_\+\+Is\+\_\+not
\item Please do not make all variables public if there is no need for them
\item Access to variable values should exclusively through Get/\+Set Methods
\end{DoxyItemize}

Methods\+:
\begin{DoxyItemize}
\item Constructor method
\begin{DoxyItemize}
\item The simplest constructor will initialize all variables to 0;
\item It is possible to overload for passing parameters
\end{DoxyItemize}
\item Camel Casing should be observed in naming all methods and they should all start with Uppercase letter
\item Only O\+NE return statement for non-\/void methods
\item Getter and Setter Methods should be named Get\mbox{[}variable\+Name\mbox{]} Set\mbox{[}variable\+Name\mbox{]} e.\+g. Get\+Number\+Of\+Photons()
\begin{DoxyItemize}
\item All Gets should return the datatype of the variable.
\item All Sets should be void
\end{DoxyItemize}
\item Memory Address methods will be named Copy\mbox{[}variable\+Name\mbox{]} e.\+g. int$\ast$ Copy\+Number\+Of\+Photons()
\begin{DoxyItemize}
\item Should always return a pointer to the variable in question.
\end{DoxyItemize}
\item Should only be relevant to data format itself. For analysis procedures please use the include and source folders N\+OT in dict/
\end{DoxyItemize}

Classes\+:
\begin{DoxyItemize}
\item All user defined classes should start with the letter I\mbox{[}Class\+Name\mbox{]}.
\item User defined classes should save the header (.h) file in the include directory and the .cpp in the src directory.
\item Any file which name starts with an I can be modified by the user
\item DO N\+OT modify A\+NY file in the dict folder except for /include/\+I\+Event.h and /src/\+I\+Event.cpp
\item Analysis classes should not need instantiation in order for them to be used
\end{DoxyItemize}

We want each stage or analysis part to be run as

Analysis\+Tool\+::\+Obtain\+Information(foo, bar)

To do so, a .h (header) file needs to be created in the include directory The definitions will go in the .cpp file in the src directory

\subsection*{E\+X\+A\+M\+P\+LE\+:}

\subsubsection*{Analysis\+Tool.\+h}


\begin{DoxyCode}
\{c++\}
        class AnalysisTool\{
            public:
                static int ObtainInformation(int foo, int bar);
        \};
\end{DoxyCode}


\subsubsection*{Analysis\+Tool.\+cpp}


\begin{DoxyCode}
\{c++\}

        int ObtainInformation(int foo, int bar)\{
            return foo*bar;
        \}
\end{DoxyCode}


By doing it this way, the dependency of different modules with each other are decreased.

A good example of this is included as \hyperlink{IPlotTools_8h_source}{I\+Plot\+Tools.\+h} and I\+Plot\+Tools.\+cpp 